%
%
%
%


%
\usepackage{lipsum}
\usepackage{amsfonts}
\usepackage{graphicx}
\usepackage{epstopdf}
\usepackage{algorithm}
\usepackage{algorithmic}
\usepackage{booktabs}
\usepackage{multirow}
\usepackage{wrapfig}
%
%
%
\ifpdf
  \DeclareGraphicsExtensions{.eps,.pdf,.png,.jpg}
\else
  \DeclareGraphicsExtensions{.eps}
\fi

\usepackage{amsmath}
\usepackage{amssymb}
\usepackage{siunitx}
\usepackage{dsfont}
\usepackage{xcolor}
\usepackage{bbm}

\usepackage{mathtools}
%

%
\newcommand{\note}[1]{{\color{blue} #1}}

%
\newcommand{\R}{\mathbb{R}}
\newcommand{\sumx}{\sum_{x\in\mathcal{X}}}
\newcommand{\tr}{\text{tr}}
\newcommand{\sumion}{\ensuremath{\sum_{i=1}^N}}
\newcommand{\sumio}[2][i]{\ensuremath{\sum_{#1=1}^#2}}
\newcommand{\avgion}[1][i]{\ensuremath{\frac{1}{N}\sum_{#1=1}^N}}
\newcommand{\prodion}{\ensuremath{\sum_{i=1}^n}}
\newcommand{\liston}[2][n]{\ensuremath{#2_{1},\dotsc,#2_{#1}}}
\newcommand{\argmin}{\ensuremath{\text{arg min }}}
\newcommand{\dom}{\mathcal{X}}

%
\newcommand{\intx}{\int_{\mathcal{X}}}
\newcommand{\pderiv}[2][ ]{\ensuremath{\frac{\partial^{#1}}{\partial #2}}}
\newcommand{\deriv}[2][ ]{\ensuremath{\frac{d^{#1}}{d #2}}}
\newcommand{\eqae}{\overset{\text{a.e.}}{=}}
\newcommand{\eqas}{\overset{\text{a.s.}}{=}}

%
\newcommand{\Pfam}{\mathcal{P}}
\newcommand{\N}{\mathbb{N}}
\newcommand{\ind}[1]{\mathbbm{1}\{#1\}}
\newcommand{\Ex}[2][ ]{\ensuremath{\mathbb{E}_{#1}\left[#2\right]}}
\newcommand{\Var}[2][ ]{\ensuremath{\text{Var}_{#1}\left(#2\right)}}
\renewcommand{\Pr}[2][ ]{\ensuremath{\mathbb{P}_{#1}\left(#2\right)}}
\newcommand{\cvgp}{\overset{p}{\rightarrow}}
\newcommand{\Xbar}{\ensuremath{\overline{X}}}
\newcommand{\Norm}[1]{\ensuremath{\mathcal{N}\left(#1\right)}}
\newcommand{\MLE}[1]{\ensuremath{\hat{#1}}_{MLE}}
\newcommand{\simiid}{\overset{\textrm{i.i.d.}}{\sim}}
\newcommand{\simind}{\overset{\textrm{ind.}}{\sim}}
\newcommand{\PCord}{\ensuremath{m}}
\newcommand{\Mmat}{\hat{M}_{\PCord,\sigma_0}}
%
\newcommand{\initialset}{\ensuremath{\mathcal{X}_0} }
\newcommand{\finalset}{\ensuremath{\mathcal{X}_1} }
\newcommand{\eventset}{\ensuremath{\mathcal{X}_e} }
\newcommand{\inputset}{\ensuremath{\mathcal{U}} }
\newcommand{\distset}{\ensuremath{\mathcal{D}} }
\newcommand{\rs}{\ensuremath{R_{[t_0,t_1]}}}
\newcommand{\ars}{\ensuremath{\hat{R}_{[t_0,t_1]}}}

%
\newcommand{\cclass}{\mathcal{C}}
\newcommand{\risk}[1][c]{\ensuremath{r(#1)}}
\newcommand{\erisk}[1][c]{\ensuremath{\hat{r}(#1)}}
\newcommand{\srisk}[1][Q]{\ensuremath{r_{#1}}}
\newcommand{\esrisk}[1][Q]{\ensuremath{\hat{r}_{#1}}}


%
\usepackage{mathtools}
\DeclarePairedDelimiter{\ceil}{\lceil}{\rceil}

\newtheorem{problem}{Problem}

%
\newcommand{\creflastconjunction}{, and~}

%
\newsiamremark{remark}{Remark}
\newsiamremark{hypothesis}{Hypothesis}
\crefname{hypothesis}{Hypothesis}{Hypotheses}
\newsiamthm{claim}{Claim}

%
\headers{Data-driven Reachability with Christoffel Functions}{A. Devonport, F.Yang, L. El Ghaoui, and M. Arcak}

%
%
\title{Data-Driven Reachability analysis and Support set Estimation with Christoffel Functions%
    \thanks{Submitted to the editors December 17, 2021. This paper is a revision and extension of a conference paper~\cite{devonport2021datadriven}.
\funding{This work is funded in part by the Air Force Office of Scientific
Research grant FA9550-21-1-0288, National Science Foundation grant ECCS-1906164,
and the Office of Naval Research grant N00014-18-1-2209.}
}}

%
\author{Alex Devonport 
    \thanks{University of California, Berkeley,
    Berkeley, CA\\ (\email{\{alex\_devonport,forestyang,elghaoui,arcak\}@berkeley.edu})}
    \and Forest Yang\footnotemark[2]
    \and Laurent El Ghaoui\footnotemark[2] 
    \and Murat Arcak\footnotemark[2]
}

\usepackage{amsopn}
\DeclareMathOperator{\diag}{diag}


%
\makeatletter
\newcommand*{\addFileDependency}[1]{%
  \typeout{(#1)}%
  \@addtofilelist{#1}%
  \IfFileExists{#1}{}{\typeout{No file #1.}}%
}
\makeatother

\newcommand*{\myexternaldocument}[1]{%
    \externaldocument{#1}%
    \addFileDependency{#1.tex}%
    \addFileDependency{#1.aux}%
}
%
%
%
%
%
